%% ---- Supplementary Note [note:dango-repro] -- OWNED BY its section agent; fill in place. ----
%% Every number here is emitted by experiments/005-kuzmin2018-tmi/scripts/dango_construction_si.py
%% (decreased zeros, lambda, dataset and split sizes; results/dango_decreased_zeros.csv and
%% results/dango_dataset_split.csv) and experiments/005-kuzmin2018-tmi/scripts/dango_string_version_sweep.py
%% (wandb pull of zhao-group/torchcell_005-kuzmin2018-tmi_dango; results/dango_string_version_{sweep,curves}.csv;
%% writes tab-dango-string-versions.tex). The figure is composed by
%% experiments/005-kuzmin2018-tmi/scripts/compose_dango_si_figures.py. Rerun those, not this file.
%% DANGO paper quotes are verbatim from the bioRxiv preprint (doi:10.1101/2020.11.26.400739), read from
%% the Zotero PDF; the citation key is not yet in references.bib (see [ref to add] below).
\clearpage
\suppnote{note:dango-repro}{constructing DANGO in TorchCell and reproducing its published result\secstatus{todo}}

DANGO predicts a trigenic interaction from six STRING evidence networks: one graph network per
network learns gene embeddings by reconstructing that network, a learned average merges them,
and a hypergraph self-attention head scores the three perturbed genes (Zhang et al.\ 2020;
\textbf{[ref to add: zhangDANGOPredictingHigherorder2020]}). It is the published model the
trigenic experiment must beat, so it was rebuilt inside TorchCell rather than run from its
released code. On the Kuzmin~2018 trigenic screen
alone~\cite{kuzminSystematicAnalysisComplex2018a} the rebuilt model reaches a validation
Pearson $r$ of $0.42$ in every one of 19 runs, whichever STRING release supplies the networks
and whichever schedule hands over from pretraining to the interaction loss
(\suppfig{fig:dango-repro}; \supptab{tab:dango-string-versions}). That plateau is close to the
value DANGO reports and cannot be equated with it, for the reasons of data, split, and protocol
listed below; what it establishes is a stable, reproducible baseline, which
\suppnoteref{note:dango-full} then retrains on the full trigenic dataset of the CGT experiment.

\notesec{Model in the manuscript's notation} DANGO reads the six channels as relation
adjacencies $A^{(k)}\in\{0,1\}^{N\times N}$ over the same $N=6607$ gene vocabulary as the CGT,
an edge wherever the channel score is nonzero (\suppnoteref{note:graphs}), but as
message-passing edges rather than attention priors (\suppfig{fig:dango-repro}a). A shared
lookup $\mathbf{E}\in\mathbb{R}^{N\times d}$, $d=64$, feeds one two-layer GraphSAGE encoder
per channel, $H^{(k)}=F^{(k)}(\mathbf{E},A^{(k)})$, and a linear head reconstructs each
channel's adjacency, $\hat A^{(k)}=H^{(k)}W_k$, under a mean squared error whose zero entries
are down-weighted by a per-channel $\lambda_k$,
\begin{equation}
\mathcal{L}_{\mathrm{rec}}=\frac{1}{6}\sum_{k=1}^{6}\frac{1}{N^{2}}\sum_{i,j}
\big(\hat A^{(k)}_{ij}-A^{(k)}_{ij}\big)^{2}\Big[\mathbf{1}\big(A^{(k)}_{ij}\neq0\big)+\lambda_k\,\mathbf{1}\big(A^{(k)}_{ij}=0\big)\Big].
\label{eq:dango-rec}
\end{equation}
A meta-embedding merges the channels per gene, $h_i=\sum_k a_{ik}h_i^{(k)}$ with
$a_{i:}=\operatorname{softmax}_k\operatorname{MLP}(h_i^{(k)})$, into a cell representation
$H$ that plays the role of Eq.~\ref{eq:encode}. There is no perturbation operator: a strain
$p=\{g_i,g_j,g_k\}$ selects the rows $h_j$, $g_j\in p$, and the Hyper-SAGNN readout scores
them from the difference of a static and an attention-derived dynamic embedding, trained with
$\mathcal{L}_{\mathrm{int}}=\operatorname{mean}\log\cosh(\hat y_{\mathrm{int}}-y_{\mathrm{int}})$.
DANGO pretrains on Eq.~\ref{eq:dango-rec} and then trains end to end; TorchCell folds both
stages into one run, $\mathcal{L}=\alpha_e\mathcal{L}_{\mathrm{rec}}+(1-\alpha_e)\mathcal{L}_{\mathrm{int}}$,
with $\alpha_e$ following one of three schedules over the first ten epochs
(\suppfig{fig:dango-repro}a, strip).

\notesec{Per-channel zero weights} The paper sets $\lambda_k$ from how many of a channel's
zeros later became edges: ``we calculated the percentage of decreased zeroes from STRING
database v9.1 to v11.0 for each network, ranging from 0.02\% (co-occurrence) to 2.42\%
(co-expression)'', and ``if the decreased percentage of zeros is greater than 1\%, we set
$\lambda=0.1$; otherwise, we set $\lambda=1.0$.'' The computation is not given, so TorchCell
fixes one: over the genes present in both releases, the pairs that were non-edges in the
older release and edges in the newer, divided by the older release's non-edge pairs. This
puts neighborhood, co-expression, and experimental above 1\% for v9.1 to v11.0 (4.60\%,
2.50\%, 1.58\%; $\lambda_k=0.1$) and fusion, co-occurrence, and database below it (0.20\%,
0.35\%, 0.33\%; $\lambda_k=1$); co-expression lands near the paper's 2.42\% and co-occurrence
does not match its 0.02\%. The same rule on v11.0 to v12.0 gives the same assignment
(\suppfig{fig:dango-repro}b). The runs below applied it only for the v9.1 networks: the
training script looked the weights up under v9.1 channel names, so the v11.0 and v12.0 runs
fell back to $\lambda_k=1$ for every channel. Found while writing this note; whether it
matters is not measured.

\notesec{Dataset, training, and result} The screen was queried from the knowledge graph
(YEPD, 30\,$^{\circ}$C, every perturbation type kept) into 91,050 triple-mutant instances over
1,400 genes, all but one with a negative interaction, and split 72,841 / 9,105 / 9,104 by
seed (train / validation / test). Runs used AdamW (learning rate $10^{-5}$, weight decay
$10^{-6}$, batch 32), $d=64$, four attention heads, and 442 to 1,000 logged epochs. The
statistic is the maximum over epochs of validation Pearson $r$, the checkpoint-selection rule
and an upward-biased order statistic whose epoch count is reported beside it; no test-split
metric was logged. Across the 19 runs it spans $0.415$ to $0.427$, with release-by-schedule
means of $0.420$ to $0.424$ and SEM at most $0.003$ where $n>1$, while training $r$ at the
selected epoch spans $0.49$ to $0.55$ (\suppfig{fig:dango-repro}c;
\supptab{tab:dango-string-versions}). Validation $r$ reaches $0.4$ within 50 epochs, peaks
between epochs 106 and 439, and then declines, to $0.32$ by epoch 1,000 in the two runs
that went that far, while training $r$ climbs to $0.55$--$0.74$ (\suppfig{fig:dango-repro}d):
the plateau is a generalization limit, not an optimization failure. No release or schedule
separates from the others beyond the run-to-run spread, although the v12.0 channels carry
three to nine times the edges of v9.1 (\suppnoteref{note:graphs}). DANGO reports about
$r=0.47$ on this screen under a random split (its Fig.~2b, read from the bar), $0.46$ at the
loosest replicate-consistency cutoff of its Fig.~2c, and a replicate ceiling for the labels:
``The Pearson correlation between the trigenic interaction scores of two individual
replicates is around 0.59''.

\notesec{What the comparison does and does not match} The architecture, the six channels at
the v9.1 release, the 1\% rule, the log-cosh loss, and the screen are the same. Five things
are not. The data are the knowledge-graph build of the screen, 91,050 instances against
``around 91,000 triplets among 1,400 genes'', not DANGO's released files; no run in TorchCell
has loaded DANGO's processed data or its folds. The evaluation is one seeded 80/10/10 split
scored on validation, not five-fold cross-validation scored on held-out folds, with no
replicate-consistency filter. The vocabulary is the 6,607-gene reference rather than the
1,395 screened genes, so the encoders reconstruct an adjacency about twenty times larger.
Pretraining and fine-tuning are one scheduled run with the optimizer settings above rather
than two stages with theirs. And the statistic is a maximum over epochs. Under these
differences the rebuilt model's $0.42$ sits about $0.05$ below the published random-split
value, a gap the split and the selection rule alone could account for; the record supports
no stronger statement in either direction.

%% Figure source: notes/assets/drawio/FigS-dango-reproduction.drawio (written by
%%   experiments/005-kuzmin2018-tmi/scripts/compose_dango_si_figures.py; panel a is draw.io cells,
%%   panel b from dango_construction_si.py, panels c and d from dango_string_version_sweep.py)
%%   -> figures/FigS-dango-reproduction.pdf (auto-exported by `make fig`/`make paper`).
\begin{figure*}[t]
\tcfig{figures/FigS-dango-reproduction.pdf}{\textbf{DANGO rebuilt in TorchCell and trained on
the Kuzmin~2018 trigenic screen.} (a)~The model in the manuscript's notation: the six STRING
channels $A^{(k)}$ are message-passing edges for one GraphSAGE encoder each, a reconstruction
head with zero weight $\lambda_k$ gives the pretraining loss, a meta-embedding merges the
channels into $H$, the perturbation selects rows of $H$ with no operator, and the Hyper-SAGNN
head scores the triple under a log-cosh loss; the strip gives the objective and the three
pretraining-to-main schedules (transition epoch 10). (b)~Percentage of decreased zeros per
channel for v9.1 to v11.0 and v11.0 to v12.0 with the resulting $\lambda_k$ printed over each
bar; dashed line, the 1\% rule; open circles, the two values the DANGO paper prints. (c)~Best
validation Pearson~$r$ (maximum over epochs) by loss schedule and STRING release; bar, mean
over runs; whisker, SEM where $n>1$; circle, one run (19 runs). (d)~Training and validation
Pearson~$r$ per epoch for the same runs; color, STRING release; line style, schedule. Sources:
\texttt{experiments/005-kuzmin2018-tmi/scripts/dango\_construction\_si.py} (b) and
\texttt{dango\_string\_version\_sweep.py} (c,~d).}{fig:dango-repro}
\end{figure*}

%% SOURCE: experiments/005-kuzmin2018-tmi/scripts/dango_string_version_sweep.py -- AUTO-GENERATED from wandb zhao-group/torchcell_005-kuzmin2018-tmi_dango; do not hand-edit.
\begin{table}[t]
\centering
\footnotesize
\caption{DANGO replication in TorchCell on the Kuzmin 2018 trigenic interactions: best
validation Pearson $r$ (maximum over epochs of the logged validation Pearson, the
checkpoint-selection rule) by STRING release and loss schedule. Mean $\pm$ SEM over $n$
runs; a single run reports its value with no whisker. Epochs is the range of epochs logged
across those runs.}
\label{tab:dango-string-versions}
\begin{tabular}{@{}l l r l l@{}}
\toprule
\textbf{STRING} & \textbf{Loss schedule} & $n$ & \textbf{Val.\ Pearson $r$} & \textbf{Epochs}\\
\midrule
v9.1 & pretrain then main & 1 & 0.420 & 1000\\
v9.1 & linear to uniform & 1 & 0.422 & 575\\
v9.1 & linear to flipped & 2 & 0.422 $\pm$ 0.001 & 500--540\\
v11.0 & pretrain then main & 1 & 0.419 & 988\\
v11.0 & linear to uniform & 2 & 0.424 $\pm$ 0.003 & 443--468\\
v11.0 & linear to flipped & 2 & 0.422 $\pm$ 0.001 & 500\\
v12.0 & pretrain then main & 4 & 0.423 $\pm$ 0.002 & 500\\
v12.0 & linear to uniform & 5 & 0.420 $\pm$ 0.001 & 500\\
v12.0 & linear to flipped & 1 & 0.422 & 500\\
\bottomrule
\end{tabular}
\end{table}

